\documentclass[twoside]{article}
\usepackage{microtype}
\usepackage{graphicx}
\usepackage{subfigure}
\usepackage{booktabs} 
\usepackage{multirow}
\usepackage{hyperref}
\usepackage{kotex}
\usepackage{amsmath}
\usepackage{amssymb}
\usepackage{graphicx}
\usepackage{subcaption} 
%\usepackage{aistats2026}
\usepackage[accepted]{aistats2026}
\usepackage[round]{natbib}
\renewcommand{\bibname}{References}
\renewcommand{\bibsection}{\subsubsection*{\bibname}}
\usepackage{amsmath}
\usepackage{amssymb}
\usepackage{mathtools}
\usepackage{algorithm}
\usepackage{algpseudocode}
\usepackage{amsthm}
\usepackage{enumitem}
\usepackage{bm}
\newcommand{\TODO}[1]{\textbf{[TODO: #1]}}
\usepackage[capitalize,noabbrev]{cleveref}
\usepackage{booktabs}
\usepackage[table]{xcolor}
\theoremstyle{plain}
\newtheorem{theorem}{Theorem}[section]
\newtheorem{proposition}[theorem]{Proposition}
\newtheorem{lemma}[theorem]{Lemma}
\newtheorem{corollary}[theorem]{Corollary}
\theoremstyle{definition}
\newtheorem{definition}[theorem]{Definition}
\newtheorem{assumption}[theorem]{Assumption}
\theoremstyle{remark}
\newtheorem{remark}[theorem]{Remark}
\newtheorem{axiom}[theorem]{Axiom}
\newtheorem{proof sketch}[theorem]{Proof Sketch}
\definecolor{darkblue}{rgb}{0.0, 0.1, 0.5}
\definecolor{blue}{RGB}{35,82,145}
\definecolor{green}{RGB}{38,126,91}
\definecolor{orange}{RGB}{184,101,29}
\begin{document}
\renewcommand{\thefootnote}{\fnsymbol{footnote}}

\twocolumn[
\aistatstitle{Adaptive Variable Selection for Continuous Treatment Effect Estimation}

\aistatsauthor{
Hyunsoo Cho\footnotemark[1] \quad
Taeseong Yoon\footnotemark[1] \quad
Insoo Jung\footnotemark[1] \quad
Heeyoung Kim
}

\aistatsaddress{
Department of Industrial and Systems Engineering, KAIST
}
]

\footnotetext[1]{Equal contribution.}
\renewcommand{\thefootnote}{\arabic{footnote}}

%\begin{abstract}
%Estimating treatment effects under continuous treatments is a fundamental problem in causal inference. Many existing continuous treatment estimators use all observed covariates to avoid omitting true confounders. However, indiscriminately incorporating all observed covariates may retain instrumental and irrelevant variables, thereby degrading estimation accuracy and obscuring the underlying causal structure. We therefore propose the Continuous Variable Selection Causal Estimation Network (CVSCEN), an end-to-end neural framework that jointly learns covariate role assignments and continuous dose--response functions without prespecifying structural priors such as the number of selected variables. A Gumbel–Softmax selector assigns each covariate to the confounder, outcome predictor, or excluded role, while one-sided treatment-dependence guidance and HSIC regularization promote role separation. A dual-stream outcome model further separates treatment-independent prognostic variation from treatment-conditioned response variation. Experiments on fully synthetic and semi-synthetic IHDP benchmarks show that CVSCEN achieves the lowest out-of-sample MISE and ADRF error among the evaluated methods while recovering role-relevant covariates with high precision and recall.
%\end{abstract}
{\color{blue}
\begin{abstract}
Estimating continuous treatment effects from observational data is a fundamental problem in causal inference.
Most continuous-treatment estimators model dose--response functions using all observed covariates as fixed inputs, potentially retaining instrumental and irrelevant variables that offer no adjustment benefit and can reduce estimation accuracy.
Existing causal variable-selection methods, however, are typically separate from effect estimation or limited to binary treatments.
To address this gap, we propose the Continuous Variable Selection Causal Estimation Network (CVSCEN), an end-to-end framework that jointly performs role-aware variable selection and continuous dose--response estimation without a predefined selection size.
For role-aware selection, CVSCEN assigns each covariate to a confounder, outcome-predictor, or excluded role using a differentiable selector guided by treatment dependence and independence regularization.
For dose--response estimation, a dual-stream outcome model uses the selected covariates to separately capture treatment-independent prognostic variation and treatment-conditioned response variation.
Experiments on fully synthetic and semi-synthetic IHDP benchmarks show that CVSCEN achieves the lowest out-of-sample MISE and ADRF error among the evaluated methods while maintaining low false discovery rates.
\end{abstract}


\section{Introduction}
Estimating causal effects from observational data is central to decision-making in medicine, econometrics, and the social sciences \citep{rosenbaum1983central,rubin2005causal}. 
Although randomized controlled trials provide a principled basis for causal estimation, practical, ethical, and financial constraints often necessitate reliance on observational data \citep{imbens2015causal}.
Under covariate adjustment, causal identification requires an appropriate set of pre-treatment covariates, but the variables needed to control confounding are rarely known a priori.
Practitioners therefore often adopt an indiscriminate covariate adjustment strategy, including all measured covariates to reduce the risk of omitting important adjustment variables \citep{vander2011new}.

%Many real-world interventions, including drug dosages and discount rates, vary continuously \citep{bica2020estimating, kazemi2024adversarially, nagalapatti2024continuous}. 
These covariate-adjustment challenges are particularly relevant in continuous-treatment settings, as many real-world interventions, including drug dosages and discount rates, vary continuously \citep{bica2020estimating, kazemi2024adversarially, nagalapatti2024continuous}.
Unlike binary settings, continuous-treatment estimation requires learning an entire dose–response function rather than comparing outcomes at only two discrete levels \citep{hirano2004propensity,imai2004causal}. 
Against this backdrop, existing approaches, however, generally focus on only one of two objectives: learning this continuous dose–response function or identifying causal covariates. Continuous-treatment estimators flexibly model outcomes across treatment levels but typically treat the observed covariates as a fixed input set, failing to explicitly distinguish their causal roles \citep{nie2021vcnet, bica2020estimating}. Conversely, variable-selection methods often rely on predictive relevance, filter covariates only implicitly, or perform selection separately from treatment effect estimation \citep{greenewald2021high, vanderschueren2025autocate}. While VSCEN \citep{ma2025simultaneous} most directly bridges this gap by jointly optimizing variable selection and treatment effect estimation end-to-end, it remains strictly limited to binary treatments and requires the number of selected variables to be specified in advance.


%Yet this strategy can be suboptimal because covariates play distinct roles.
Treating all covariates alike can be suboptimal because they play distinct causal roles.
Variables needed for confounding adjustment must be retained, while pure outcome predictors may improve outcome modeling and estimation efficiency.
By contrast, once a sufficient adjustment set is available, including treatment-only and irrelevant variables primarily inflates finite-sample variance and high-dimensional noise without improving confounding control. 
In particular, indiscriminately adjusting for instrumental variables not only amplifies the bias of unobserved confounders \citep{pearl2011invited,ding2017instrumental} but is also especially detrimental in continuous-treatment settings; incorporating instruments forces the outcome model to condition on strong, non-prognostic treatment-assignment signals, which exacerbates variance inflation in regions of limited overlap and severely destabilizes the estimation of the continuous dose-response curve. 
Taken together, these challenges motivate causal role-aware variable selection, which retains variables essential for adjustment or outcome modeling while actively filtering out those that degrade estimation accuracy \citep{pearl2014confounding,ma2025simultaneous}.
Beyond selecting role-relevant covariates, continuous dose--response estimation also calls for an outcome model that distinguishes treatment-independent prognostic variation from treatment-conditioned response variation.

%Directly extending such binary variable-selection methods to continuous domains is nontrivial. In binary settings, treatment effects are estimated as discrete shifts, whereas learning a continuous dose-response curve requires a neural network to model complex, non-linear interactions across the entire dosage continuum. While failing to filter instruments destabilizes estimation as previously noted,  this continuous functional mapping also makes the model highly vulnerable to entangling the representations of treatment signals and prognostic covariates. If continuous treatment signals strongly correlate with certain prognostic covariates, the outcome model can easily misattribute baseline prognostic variations to the treatment effect, thereby severely biasing the estimated dose-response curve. Thus, there remains a critical need for an end-to-end framework that adaptively assigns covariates to causal roles without predefined cardinality constraints while robustly learning the continuous dose–response function. 

To address these two requirements, we propose the \textit{Continuous Variable Selection Causal Estimation Network} (CVSCEN), an end-to-end framework that jointly performs causal role-aware variable selection and continuous dose--response estimation.
For role-aware variable selection, CVSCEN uses a dynamic Gumbel--Softmax selector to assign each covariate to one of three roles---confounder, outcome predictor, or excluded.
Because confounders and outcome predictors can both improve factual outcome prediction, one-sided treatment-dependence guidance suppresses weakly treatment-associated variables from the confounder role without directly rewarding strong association.
An HSIC-based penalty further discourages treatment-assignment information from entering the outcome-predictor representation, while delayed role-specific sparsity regularization encourages a compact selected set without prespecifying its cardinality.
For continuous dose--response estimation, CVSCEN employs a dual-stream outcome model that decomposes prediction into treatment-independent baseline and treatment-conditioned response components.
Multi-frequency treatment encoding and treatment-conditioned affine modulation are applied to the response stream to preserve the scalar treatment signal and model nonlinear variation over the treatment domain.
The selector and outcome model are trained jointly under a combined objective comprising factual outcome loss, HSIC-based independence regularization, and delayed role-specific sparsity penalties.
We further provide a finite-sample theoretical analysis showing that the ADRF estimation error decomposes into residual confounding from selected covariates, factual outcome-model learning error, and empirical averaging error. 
The residual confounding term is controlled by the product of residual treatment-assignment and outcome sensitivities of omitted covariates, formalizing that omitted variables are harmful mainly when they are both treatment-associated and outcome-relevant. 
Thus, the analysis supports CVSCEN's role-aware design without requiring exact recovery of every covariate role.


We evaluate CVSCEN on a fully synthetic benchmark with controlled causal roles and a semi-synthetic IHDP benchmark with real-world covariates.
Across both settings, CVSCEN achieves the lowest MISE and ADRF error among the evaluated methods while retaining role-relevant covariates and largely excluding treatment-only and irrelevant variables.

{\color{blue}
\section{Related Work}
\subsection{Continuous Treatment Effect Estimation}
Continuous treatment effect estimation is commonly studied through the generalized propensity score and related potential-outcome frameworks \citep{hirano2004propensity, imai2004causal}. 
These approaches provide a principled foundation for identifying dose-response functions, but often require accurate modeling of the treatment assignment mechanism and can be sensitive to limited overlap in observational data. 

Recent deep learning methods model complex nonlinear dose-response relationships using flexible neural architectures. 
SCIGAN uses a generative adversarial framework to estimate counterfactual outcomes under continuous interventions \citep{bica2020estimating}, whereas VCNet employs varying-coefficient neural networks combined with targeted regularization to estimate individualized dose-response curves \citep{nie2021vcnet}. 
DRNet \citep{schwab2020learning} partitions the continuous treatment space into dosage intervals and learns separate prediction heads for counterfactual estimation. 
Other methods improve robustness through adversarially balanced representations \citep{kazemi2024adversarially}, gradient interpolation techniques \citep{nagalapatti2024continuous}, or proximal causal learning \citep{wu2023doubly}. 
Despite these advances, most continuous-treatment estimators treat the observed covariate set as fixed and do not distinguish variables needed for confounding adjustment or outcome prediction from treatment-only or irrelevant variables, potentially reducing estimation stability in high-dimensional settings.


\subsection{Variable Selection for Causal Inference}
Variable selection has been widely studied in causal inference to improve sample efficiency and reduce the influence of irrelevant or harmful covariates. 
Classical regularization methods, such as the Lasso, can remove weakly predictive variables through coefficient shrinkage \citep{tibshirani1996regression}. 
However, prediction-oriented selection does not account for causal roles and therefore cannot structurally distinguish confounders, outcome predictors, and treatment-only variables such as instruments \citep{greenewald2021high}. 

Robust causal estimators, including Double/Debiased Machine Learning \citep{chernozhukov2018double} and Generalized Random Forests (GRF) \citep{athey2019generalized}, use flexible nuisance models to provide reliable treatment effect estimates in high-dimensional settings, with recent extensions to continuous treatments \citep{zenati2025double}. 
However, these methods are primarily designed for robust estimation rather than explicit covariate-level causal role assignment.
Representation-based methods and automated pipelines can reduce the influence of nuisance covariates: AutoCATE \citep{vanderschueren2025autocate} automates feature and model selection, while IBEX \citep{van2026continuous} suppresses non-causal information through an information bottleneck objective. 
However, they either separate variable selection from treatment effect estimation or provide only implicit feature filtering without interpretable covariate-level role assignments.


More closely related to our work, VSCEN \citep{ma2025simultaneous} jointly estimates treatment effects and selects causal variables end-to-end using Concrete relaxations \citep{jang2017categorical,maddison2017concrete}.
However, it is limited to binary treatments and requires a predefined selection cardinality. 
CVSCEN extends this framework to continuous treatment effect estimation by adaptively learning role-specific masks for adjustment and outcome modeling, without fixing the number of selected variables. 
}

\section{Background}
\subsection{Continuous Treatment Effect Estimation}
We consider causal effect estimation with continuous treatment under the potential outcome framework \citep{rubin2005causal, imbens2015causal}. 
For each unit $i$, let $A_i \in \mathcal{A}$ denote the observed continuous treatment, where $\mathcal{A}$ is the treatment domain, and let $Y_i^a$ denote the potential outcome under treatment $A_i = a$.

Unlike binary treatment settings, where the causal effect is often summarized by a contrast between two treatment conditions, continuous treatment effect estimation aims to infer how potential outcomes vary over a continuum of treatment levels.
This relationship is described by a dose--response function, which maps each treatment level $a$ to the expected potential outcome \citep{hirano2004propensity, imai2004causal}.
In this work, we focus on the average dose--response functions, which summarize the expected potential outcome at each treatment level.

At the population level, the average dose--response function (ADRF) is defined as
\begin{equation}
    \mu(a):=\mathbb{E}[Y^a]
\end{equation}
which represents the expected potential outcome if all units in the population were assigned treatment level $a$.

To capture treatment effect heterogeneity across covariates, we define the conditional average dose--response function (CADRF) as
\begin{equation}
    \mu(a,\mathbf{x}):=\mathbb{E}[Y^a\mid\mathbf{X}=\mathbf{x}]
\end{equation}
which represents the expected potential outcome at treatment level $a$ for units with covariates \(\mathbf{X}=\mathbf{x}\).

The ADRF can be obtained by marginalizing the CADRF over the covariate distribution:
\begin{equation}
    \mu(a) = \mathbb{E}_{\mathbf{X}}[\mu(a,\mathbf{X})]
\end{equation}

Thus, estimating the CADRF allows us to recover a conditional dose--response curve for each covariate profile, while the ADRF summarizes the corresponding population-level response.

\subsection{Identification from Observational Data}
Identifying dose--response functions from observational data requires assumptions that connect the observed data distribution to the potential outcome distribution
\citep{rosenbaum1983central, hirano2004propensity, imai2004causal}.

\begin{itemize}
    \item \textbf{Consistency.} Treatment level $A_i=a$ implies $Y_i=Y_i^a$ for unit $i$.
    \item \textbf{Weak unconfoundedness.} For every treatment level $a \in \mathcal{A}$ and unit $i$, $Y_i^a \perp\!\!\!\perp A_i \mid \mathbf{X}_i$.
    \item \textbf{Positivity.} For all covariate values $\mathbf{x}\in\mathbf{X}$ and treatment levels $a$, 
    $f_{A\mid \mathbf{X}}(a\mid \mathbf{x})>0$.
\end{itemize}

In addition to these standard identification assumptions, we assume that the observed pre-treatment covariates admit a causal-role partition:
\begin{equation}
    \mathbf{X} = \mathbf{X}_C \cup \mathbf{X}_P \cup \mathbf{X}_I \cup \mathbf{X}_N,
\end{equation}
where $\mathbf{X}_C$, $\mathbf{X}_P$, $\mathbf{X}_I$, and $\mathbf{X}_N$ denote confounders, outcome predictors, instrumental variables, and irrelevant variables, respectively.

Confounders affect both the treatment and the outcome, and therefore should be adjusted to remove confounding bias. Outcome predictors affect the outcome, but not the treatment, and including them can improve outcome modeling and finite-sample efficiency. %\cite{}

Let $\textbf{S} = \mathbf{X}_C \cup \mathbf{X}_P$ denote the role-relevant covariate set.
We assume that $\textbf{S}$ is \(c\)-equivalent to the full observed covariate set \(\mathbf{X}\), meaning that adjusting for $\textbf{S}$ provides the same confounding adjustment effect as adjusting for all observed covariates \citep{pearl2014confounding, greenewald2021high}.

In contrast, instrumental variables $\textbf{X}_I$, which affect the treatment but do not directly affect the outcome, are not required for covariate adjustment and may adversely affect finite sample estimation by increasing the variance of treatment or propensity score models.
Irrelevant variables $\mathbf{X}_N$, which affect neither the treatment nor the outcome, do not provide causal or predictive information and may further increase noise, dimensionality, and estimation variance. %\cite{}

\subsection{End-to-End Variable Selection via Concrete Relaxation}
Variable selection is inherently discrete and non-differentiable, thus making it difficult to optimize jointly with a neural network.
A common strategy is to replace discrete selection variables with differentiable stochastic relaxations, such as the Concrete or Gumbel--Softmax distribution \citep{maddison2017concrete,jang2017categorical}.

Let \(\boldsymbol{\pi}=(\pi_1,\ldots,\pi_K)\) denote the class probabilities of a categorical random variable with \(K\) categories, where \(\pi_k \geq 0\) and \(\sum_{k=1}^{K}\pi_k=1\). 
A relaxed categorical sample \(\boldsymbol{y}=(y_1,\ldots,y_K)\) is then defined as
\begin{equation}
    y_k =
    \frac{
        \exp\left((\log \pi_k + g_k)/\tau\right)
    }{
        \sum_{m=1}^{K}
        \exp\left((\log \pi_m + g_m)/\tau\right)
    },
    k=1,\ldots,K,
\end{equation}
where \(g_1,\ldots,g_K\) are independent standard Gumbel random variables and \(\tau>0\) denotes the temperature.

The resulting vector \(\boldsymbol{y}\) lies on the probability simplex and can be used as a differentiable approximation to a one-hot categorical sample. 
When \(\tau\) is large, the relaxation is smoother, and probability mass is spread across categories. 
As \(\tau\) approaches zero, the relaxed sample becomes increasingly discrete and converges toward a one-hot vector with $P(y_k=1)=\pi_k/\sum_m\pi_m$.

This relaxation enables discrete variable selection decisions to be learned end-to-end within a neural architecture. 
For any positive temperature \(\tau\), these relaxed assignments act as differentiable masks, allowing gradients from the downstream treatment effect estimation objective to update the selection parameters. With temperature annealing, the selector initially explores multiple variable assignments and gradually concentrates on the variables most useful for estimating potential outcomes. 
After training, the resulting hard assignments provide an interpretable categorization of covariates and are used to construct the final adjustment set. 


\section{Continuous Variable Selection Causal Estimation Network}
\label{sec:cvscen}
\subsection{Overview}
CVSCEN is an end-to-end neural framework for continuous treatment effect estimation with causal role-aware variable selection. 
Given an observational dataset $\mathcal{D}=\{(\mathbf{x}_i,a_i,y_i)\}_{i=1}^{n}$, where \(\mathbf{x}_i\in\mathbb{R}^{d}\) denotes pre-treatment covariates, \(a_i\in\mathcal{A}\) denotes a continuous treatment, and \(y_i\in\mathbb{R}\) denotes the observed outcome, CVSCEN jointly learns covariate role assignments and a treatment-conditioned outcome model. 

Each covariate is assigned to one of three roles $\mathcal{R}=\{0,C,P\}$, corresponding to the excluded, confounder, and outcome-predictor roles.
Covariates assigned to $C$ or $P$ form the role-relevant set $\mathbf{S} = \mathbf{X}_C \cup \mathbf{X}_P$, which is used by the downstream dose-response estimator, whereas covariates assigned to $0$ are excluded.

Using the selected covariates $\mathbf{S}$, CVSCEN parameterizes the CADRF as 
\[
    \widehat{\mu}_{\Theta}(a,\mathbf{s}_i)
    \approx
    \mu(a,\mathbf{s}_i)
    =
    \mathbb{E}[Y^a \mid \mathbf{S}=\mathbf{s}_i]
\]
where \(\mathbf{s}_i\) denotes the selected covariates for unit \(i\), and \(\Theta\) collects the parameters of the variable-selection and outcome-prediction modules.

The corresponding ADRF estimate and the factual outcome prediction are:
\[
    \widehat{\mu}(a)
    =
    \frac{1}{n}
    \sum_{i=1}^n
    \widehat{\mu}_{\Theta}(a,\mathbf{s}_i), \quad \widehat{y}_i
    =
    \widehat{\mu}_{\Theta}(a_i,\mathbf{s}_i)
\]
CVSCEN comprises three components: (i) a variable selector based on the Gumbel-Softmax (Concrete) relaxation, together with one-sided treatment-dependence guidance (\cref{sec:dynamic_variable_selection}), (ii) an HSIC-based regularizer that discourages the outcome-predictor representation from retaining treatment-assignment information (\cref{sec:dynamic_variable_selection}), and
(iii) a dual-stream, treatment-conditioned outcome model operating on the selected covariates(\cref{sec:outcome_prediction}).
All components are optimized jointly using factual outcome loss, an HSIC penalty, and delayed role-specific sparsity regularization (\cref{sec:end_to_end_training}).


\subsection{Dynamic Variable Selection Layer}
\label{sec:dynamic_variable_selection}
\textbf{Role Assignment.} 
For each pre-treatment covariate \(X_j\), CVSCEN represents each covariate's role assignment over $\mathcal{R}=\{0,C,P\}$ using a learnable logit vector $\boldsymbol{\ell}_j=(\ell_{j,0},\ell_{j,C},\ell_{j,P})\in\mathbb{R}^{3}$,  where \(\ell_{j,r}\) is the assignment score of covariate \(X_j\) to role \(r\in\mathcal{R}\).

We obtain differentiable stochastic assignments using the Gumbel-Softmax relaxation \citep{jang2017categorical, maddison2017concrete}. 
At epoch \(e\), one-sided dependence guidance modifies only the confounder logit
\begin{equation}
\label{Eq7}
    \widetilde{\ell}_{j,r}^{(e)} = \ell_{j,r} + \mathbf{1}\{\tau^{(e)} \leq \tau_g\} \mathbf{1}\{r=C\} \lambda^{(e)} \widetilde{\rho}_j,
\end{equation}
where $\tau^{(e)}$ is current temperature, $\tau_g$ is the guidance activation threshold, $\lambda^{(e)}\geq 0$ controls the strength of the guidance, and the $\widetilde{\rho}_j$ denotes the one-sided dependence score defined below in \cref{eq:one_sided_score}.


With independent standard Gumbel variables $g_{j,r}^{(e)}$, the relaxed assignment probability is 
\begin{equation}
    z_{j,r}^{(e)}=\frac{\exp\left[(\widetilde{\ell}_{j,r}^{(e)} + g_{j,r}^{(e)})/\tau^{(e)}\right]}{\sum_{r'\in\mathcal{R}}\exp\left[(\widetilde{\ell}_{j,r'}^{(e)} + g_{j,r'}^{(e)})/\tau^{(e)}\right]}, \quad r\in\mathcal{R}.
\end{equation}
The vector $\boldsymbol{z}_j^{(e)}=(z_{j,0}^{(e)},z_{j,C}^{(e)},z_{j,P}^{(e)})$
lies on the probability simplex and provides a differentiable approximation to a discrete role assignment.
For $r \in \{C,P\}$, CVSCEN defines the role-specific mask $m^{r,(e)}=(z_{1,r}^{(e)},\ldots,z_{d,r}^{(e)})$ and the corresponding masked covariates $\mathbf{x}_{i}^{r}=m^{r,(e)}\odot \mathbf{x}_{i}$, where \(\odot\) denotes element-wise multiplication. 
The downstream outcome model receives separate soft representations of the selected confounders and outcome predictors, rather than the full covariate vector.


\textbf{Dependence-Based Guidance for Confounder Selection.}
Because both confounders and outcome predictors can improve factual outcome prediction, the outcome loss alone may not reliably separate their roles. 
CVSCEN therefore uses a covariate-treatment dependence as a soft prior for confounder assignment. 

For each covariate \(X_j\), let \(\mathbf{x}_{:j}=(x_{1,j},\ldots,x_{n,j})^\top\) denote its observed sample vector, and let
\(\mathbf{a}=(a_1,\ldots,a_n)^\top\) denote the observed treatment vector.  
We define its empirical treatment-dependence score as $\rho_j:=\widehat{\operatorname{HSIC}}\left(\mathbf{x}_{:j},\mathbf{a}\right)$, where $\operatorname{HSIC}$ denotes the Hilbert Schmidt Independence Criterion \citep{gretton2005measuring, gretton2007kernel}. 
These scores are standardized across covariates and transformed into a bounded one-sided guidance score:
\begin{equation}
\label{eq:one_sided_score}
    \widetilde{\rho}_j = \min \left\{{0,\tanh \left(\frac{\rho_j-\mu_{\rho}}{\sigma_{\rho}+\epsilon} \right)} \right \},
\end{equation}
where \(\mu_{\rho}\) and \(\sigma_{\rho}\) are the empirical mean and standard deviation \(\{\rho_j\}_{j=1}^{d}\), and \(\epsilon>0\) ensures numerical stability.


Because \(\widetilde{\rho}_j\leq 0\), below-average treatment dependence reduces the confounder logit, whereas above-average dependence is not directly rewarded. 
We anneal the temperature according to $\tau^{(e)} = \max({\tau_{\min}, \tau_{\mathrm{init}}\exp(-\gamma e)})$, and activate the guidance only when $\tau^{(e)} \leq \tau_g$, allowing the selector to explore role assignments before incorporating dependence information. 
This asymmetric design avoids promoting treatment-only variables such as instruments, which may be strongly associated with treatment despite being unnecessary for standard covariate adjustment and potentially harmful in finite samples. The guidance therefore acts only as a soft structural prior, while the final role assignments are learned through factual outcome fitting, HSIC regularization, and role-specific sparsity penalties.

\subsection{Outcome Prediction Layer}
\label{sec:outcome_prediction}
The outcome prediction layer is trained to estimate factual outcomes using the observed treatment and the selected role-specific covariates.
Let $\mathbf{s}_i = (\mathbf{x}_i^C,\mathbf{x}_i^P)$ denote the selected covariates used by the outcome model.
The prediction layer defines a treatment-conditioned function $F_{\Theta}(a,\mathbf{s}_i)$, and the factual prediction is obtained by evaluating this function at the observed treatment:
$\widehat{y}_i = F_{\Theta}(a_i, \mathbf{s}_i)$.
Although the layer is trained through factual outcome prediction, the learned function can later be evaluated at arbitrary treatment levels to induce an estimated CADRF.

\textbf{Treatment-Conditioned Affine Modulation.}
Continuous treatment effect estimation requires the prediction model to capture how outcomes vary with the treatment level.
Simply concatenating a scalar treatment with a high-dimensional covariate representation may weaken the treatment signal.
To address this issue, CVSCEN uses treatment-conditioned affine modulation to inject treatment information into the treatment-response component of the outcome model \citep{perez2018film}.

For a normalized treatment value $a\in[0,1]$, we first construct a multi-frequency treatment encoding with $L$ frequency bands \citep{mildenhall2020nerf}:
\begin{equation}
    \phi(a) = [a, \{\sin(2^{\ell}\pi a),\cos(2^{\ell}\pi a)\}_{\ell=0}^{L-1}]^{\top} \in \mathbb{R}^{1+2L}.
\end{equation}
This positional encoding provides the model with treatment features at multiple frequency scales, allowing it to represent nonlinear variation over the treatment domain \citep{rahaman2019spectral, tancik2020fourier, mildenhall2020nerf}.

Feature-wise scale and shift parameters are generated as $\gamma(a) = g_{\gamma}(\phi(a)), \ \delta(a) = g_{\delta}(\phi(a)).$
For latent representation $h$, the treatment-conditioned affine modulation is applied as
\begin{equation}
    \operatorname{Mod}(h,a) = h \odot (1+\gamma(a)) + \delta(a).
\end{equation}
The residual scaling term $1+\gamma(a)$ keeps the transformation close to the identity when $\gamma(a)$ and $\delta(a)$ are near zero, which helps stabilize training.

\textbf{Dual-Stream Outcome Prediction.}
CVSCEN decomposes the prediction of factual outcomes into a treatment-independent baseline component and a treatment-dependent response component.
The selected confounders and outcome predictors are first projected into latent representations: $h_i^C = \Phi_C(\mathbf{x}_i^C)$, and $h_i^P = \Phi_P(\mathbf{x}_i^P)$, where $\Phi_C$ and $\Phi_P$ are learnable projection layers.
These representations are then mapped into two stream-specific features, $h_{i, \mathrm{base}} = q_{\mathrm{base}}(h_i^C,h_i^P)$, and $h_{i,\mathrm{resp}} = q_{\mathrm{resp}}(h_i^C,h_i^P)$,
where $q_{\mathrm{base}}$ and $q_{\mathrm{resp}}$ are nonlinear adapters for the baseline and treatment-response streams, respectively. 

The baseline stream predicts the treatment-independent prognostic component, $y_{\mathrm{base},i} = f_{\mathrm{base}}(h_{i,\mathrm{base}})$. 
Because the observed treatment is not provided to this stream, it captures prognostic variation explained by the selected covariates alone.

In contrast, the treatment-response stream models treatment-dependent variation by applying treatment-conditioned affine modulation to its representation:
\begin{equation}
    y_{\mathrm{resp},i}(a_i) = f_{\mathrm{resp}}(\operatorname{Mod(h_{i,\mathrm{resp}},a_i)})
\end{equation}
The factual outcome prediction is then obtained by summing the two components:
\begin{equation}
    \widehat{y}_i = y_{\mathrm{base},i} + y_{\mathrm{resp},i}(a_i).
\end{equation}


During training, the model is evaluated at the observed treatment $a_{i}$, yielding the factual prediction $\widehat{y}_{i} =F_{\Theta}(a_{i},\mathbf{s}_i)$. 
After training, it can be evaluated at any treatment level $a\in\mathcal{A}$, providing the estimated CADRF:
\begin{equation}
    \widehat{\mu}_{\Theta}(a,\mathbf{s}_i) := F_{\Theta}(a,\mathbf{s}_i).
\end{equation}

\subsection{End-to-End Training of CVSCEN}
\label{sec:end_to_end_training}
CVSCEN is trained end-to-end to jointly learn role-aware variable selection and factual outcome. 
The training objective combines a factual outcome prediction loss, an HSIC independence penalty, and a sparsity regularization term:
\begin{equation}
    \mathcal{L} = \mathcal{L}_{Y} + \lambda_{\mathrm{hsic}}\mathcal{L}_{\mathrm{hsic}} + w^{(e)}\mathcal{L}_{\mathrm{reg}},
\end{equation}
where $\lambda_{\mathrm{hsic}}\geq 0$ controls the strength of the HSIC penalty and $w^{(e)}$ is an epoch-dependent coefficient that controls the contribution of the sparsity regularization term during training. 

\textbf{Factual Outcome Loss.}
The factual outcome loss fits the observed outcomes using the observed treatment values.
To reduce sensitivity to outliers, we employ the Huber loss \citep{huber1992robust}:
\begin{equation}
    \mathcal{L}_{Y} = \frac{1}{n}\sum_{i=1}^{n}\ell_{\delta}(y_i-\widehat{y}_i),
\end{equation}
where the Huber loss is given by \(\ell_{\delta}(z) = \frac{1}{2}z^2\) if \(|z| \leq \delta\), and \(\delta(|z| - \frac{1}{2}\delta)\) otherwise. In our implementation, we set $\delta=1.0$.

\textbf{HSIC Independence Penalty.}
To promote role separation between confounders and outcome predictors, CVSCEN penalizes statistical dependence between the masked prognostic covariates and the observed treatment:
\begin{equation}
    \mathcal{L}_{\mathrm{hsic}} = \widehat{\operatorname{HSIC}}(\{x_i^P\}_{i=1}^n,\{a_i\}_{i=1}^n),
\end{equation}
where $x_i^P = x_i \odot m_i^P$ denotes the raw covariates multiplied by the outcome predictor mask.
Minimizing this term computationally enforces the variable selection layer to drop covariates correlated with the treatment from the prognostic set, promoting a clearer distinction between prognostic and treatment-related information. 

\textbf{Delayed Sparsity Regularization.}
Applying sparsity regularization from the beginning of training may remove informative covariates before their predictive utility is learned. 
Therefore, CVSCEN activates the sparsity penalty only after the Gumbel--Softmax temperature has decreased below a threshold. 
We define $w^{(e)} = \mathbf{1}(\tau^{(e)} \leq \tau_s)$,
where $\tau_s$ is the sparsity activation threshold, set to $1.0$ by default.

The sparsity regularization term penalizes the expected number of variables assigned to the confounder and outcome-predictor roles:
\begin{equation}
    \mathcal{L}_{\mathrm{reg}} = \lambda_{\mathrm{reg},C}\sum_{j=1}^{d}z_{j,C}^{e} + \lambda_{\mathrm{reg},P}\sum_{j=1}^{d}z_{j,P}^{e},
\end{equation}
where $z_{j,C}^{e}$ and $z_{j,P}^{e}$ are the relaxed assignment probabilities of covariate $X_j$ to the confounder and outcome-predictor roles, respectively.
The weights $\lambda_{\mathrm{reg},C}$ and $\lambda_{\mathrm{reg},P}$ control the sparsity levels of the two selected variable sets independently.
This term encourages the model to retain a compact set of role-relevant covariates while assigning uninformative variables to the excluded role.

\textbf{Post-Training Discretization and Estimation.}
After training, each covariate is assigned to the role with the largest learned logit, $\widehat{r}_j=\arg\max_{r\in\mathcal{R}}\ell_{j,r}$. 
The corresponding hard masks are defined as $\widehat{m}_{C,j} = \mathbf{1}\{\widehat{r}_j=C\}$ and $\widehat{m}_{P,j} = \mathbf{1}\{\widehat{r}_j=P\}$.
Let $\widehat{\mathbf{m}}_C$ and $\widehat{\mathbf{m}}_P$ denote the resulting mask vectors. 
The final selected covariate representation for unit $i$ is then given by $\widehat{s}_i = (\widehat{\mathbf{m}}_C\odot \mathbf{x}_i, \widehat{\mathbf{m}}_P \odot \mathbf{x}_i)$.
Then, the CADRF estimator is
\begin{equation}
    \widehat{\mu}_{\Theta}(a, \mathbf{\widehat{s}}_i) := F_{\Theta}(a, \widehat{\mathbf{s}}_i),
\end{equation}
and the ADRF is estimated by
\begin{equation}
    \widehat{\mu}(a) = \frac{1}{n}\sum_{i=1}^{n}\widehat{\mu}_{\Theta}(a, \widehat{\mathbf{s}}_i).
\end{equation}
}
{\color{blue}
\section{Theoretical Analysis}
We provide a finite-sample error bound for the ADRF estimator induced by CVSCEN.
Our bound decomposes the ADRF error into three terms.
First, let $\Delta_A(\hat S)$ and $\Delta_Y(\hat S)$ bound the residual treatment-assignment and outcome sensitivities of omitted covariates after selecting $\hat S$.
With selected-covariate positivity constant $c_A>0$, define $\mathcal E_{\rm sel}:=\frac{1}{c_A}\Delta_A(\hat S)\Delta_Y(\hat S)$.
This term captures residual confounding due to variable selection.

Second, because CVSCEN is trained with the Huber factual loss, let $\mathcal L_n^H$ be the corresponding Huber-loss class induced by the outcome model and selection maps, uniformly bounded by $B_H$. 
Here, $\mathfrak R_n(\cdot)$ denotes empirical Rademacher complexity, which controls the uniform concentration of a function class. 
Define
\[
    \gamma_H(n,\delta):=2\mathfrak R_n(\mathcal L_n^H)+B_H\sqrt{\frac{\log(4/\delta)}{2n}}.
\]
Let $\mathcal B_{\rm app}$ and $\mathcal B_{\rm opt}$ denote approximation and optimization error bounds measured in Huber risk.
For example, under smoothness of the selected-covariate outcome regression, $\mathcal B_{\rm app}$ can be taken as a standard neural approximation bound, while $\mathcal B_{\rm opt}$ can be set to the empirical optimization tolerance.
Under Huber-risk consistency, local curvature, and a treatment-domain transfer condition, the prediction error is bounded by
\[
    \mathcal E_{\rm pred}:=C_{\rm tr}\sqrt{\mathcal B_{\rm app}+2\gamma_H(n,\delta)+\mathcal B_{\rm opt}}.
\]

Third, empirical averaging introduces an additional error because the population covariate distribution is replaced by the empirical distribution. 
Let $\Gamma_{\rm avg}(n,\delta)$ be a high-probability uniform concentration bound satisfying
\[
    \sup_{a\in\mathcal A}\left|\frac{1}{n}\sum_{i=1}^{n}\hat\mu_\Theta(a,\hat s_i)-\mathbb{E}[\hat\mu_\Theta(a,\hat S)]\right|\leq\Gamma_{\rm avg}(n,\delta)
\]
with probability at least $1-\delta$. 
We define $\mathcal E_{\rm avg}:=\Gamma_{\rm avg}(n,\delta)$.
Finally, set $\varepsilon_{\rm cv} := \mathcal E_{\rm sel} + \mathcal E_{\rm pred} + \mathcal E_{\rm avg}$.

\begin{theorem}[Finite-Sample ADRF Error Bound]
\label{thm:finite_adrf_bound} 
Under selected-covariate positivity, residual sensitivity control, and the prediction and averaging bounds defined above, with probability at least $1-\delta$,
\[    
    \sup_{a\in\mathcal A} \left| \hat{\mu}(a)-\mu(a) \right| \leq \varepsilon_{\rm cv}. 
\] 
\end{theorem} 
Theorem~\ref{thm:finite_adrf_bound} decomposes the ADRF estimation error into residual confounding from variable selection, factual outcome-model error controlled through the Huber risk, and empirical averaging error.
The selection term depends on the product $\Delta_A(\hat S)\Delta_Y(\hat S)$, reflecting that omitted covariates induce bias primarily when they are both treatment-associated and outcome-relevant.
Thus, the result does not require exact recovery of every causal role; it requires the selected covariates to leave small residual confounding.
The proof and concrete concentration bounds for $\gamma_H(n,\delta)$ and $\Gamma_{\rm avg}(n,\delta)$ are provided in Appendix~\ref{Proofs for the Theorems}.

}

{\color{orange}
\section{Experiments}
\subsection{Overview}
We organize the experiments around three research questions:
\begin{itemize}[leftmargin=*, itemsep=1pt, topsep=2pt, parsep=0pt]
    \item \textbf{RQ1 (Treatment Effect Estimation).} How accurately does CVSCEN estimate conditional and population-level dose-response functions?

    \item \textbf{RQ2 (Variable Selection).} Does CVSCEN retain covariates relevant to confounding adjustment or outcome modeling while excluding instrumental and irrelevant variables? 

    \item \textbf{RQ3 (Mechanism Analysis).} How do dependence-based guidance and HSIC regularization contribute to role-aware variable selection and dose-response estimation? 
\end{itemize}

\textbf{Experimental Setting}
We evaluate CVSCEN on two benchmarks:
(i) a fully synthetic dataset with controlled covariate roles and known dose-response functions (\cref{Exp: fully synthetic}), and 
(ii) a semi-synthetic dataset that combines real-world covariates with synthetically generated treatments and outcomes (\cref{Exp: semi-synthetic}).
Both settings provide ground-truth dose-response functions and causal roles, enabling quantitative evaluation of treatment effect estimation and variable selection, which cannot be directly assessed using fully observational data. 

\textbf{Baselines.}
We compare CVSCEN with representative continuous treatment effect estimators, including VCNet \citep{nie2021vcnet}, DRNet \citep{schwab2020learning}, SCIGAN \citep{bica2020estimating}, GIKS \citep{nagalapatti2024continuous}, IBEX \citep{van2026continuous}, DDMLCT \citep{colangelo2025double}, and ACFR \citep{kazemi2024adversarially}.

\textbf{Implementation Details.}
CVSCEN parameterizes role assignments using a \(d\times3\) logit matrix and anneals the Gumbel--Softmax temperature from \(10.0\) to \(0.1\).
The dual-stream outcome model uses residual MLP blocks with LayerNorm and SiLU activations, NeRF-style treatment encoding, and treatment-conditioned affine modulation.
Using a train--validation--test split, we tune all methods on the validation set via Optuna and report test-set performance.
Unlike general deep learning, hyperparameter tuning in causal inference tasks is uniquely challenging due to the absence of ground-truth counterfactuals. However, since our benchmarks provide the true dose-response functions, we directly use MISE as the tuning metric for all methods.


\textbf{Evaluation Metrics.}
Let \(\widehat{\mu}_i(a)\) denote the dose--response estimate produced by a given method for evaluation unit \(i\); for CVSCEN,
\(\widehat{\mu}_i(a)=\widehat{\mu}_{\Theta}(a,\widehat{\mathbf{s}}_i)\).
Using an equally spaced treatment grid \(\mathcal{G}\) over the 5th--95th percentiles of the observed treatment distribution, we report
\vspace{-1.5em}
\[
\begin{aligned}
\mathrm{MISE}
&=
\frac{1}{N_{\mathrm{ev}}|\mathcal{G}|}
\sum_{i=1}^{N_{\mathrm{ev}}}
\sum_{a\in\mathcal{G}}
\left[
\widehat{\mu}_i(a)-\mu(a,\mathbf{x}_i)
\right]^2,
\\
\epsilon_{\mathrm{ADRF}}
&=
\frac{1}{|\mathcal{G}|}
\sum_{a\in\mathcal{G}}
\left|
\widehat{\mu}(a)-\mu(a)
\right|,
\end{aligned}
\]
where
\(\widehat{\mu}(a)=
N_{\mathrm{ev}}^{-1}
\sum_{i=1}^{N_{\mathrm{ev}}}
\widehat{\mu}_i(a)\).
For each role \(r\in\{C,P,C\cup P\}\), we additionally report false discovery rate (FDR) and true positive rate (TPR).
%their formal definitions are provided in the Appendix.

\textbf{Evaluation Range.} 
To avoid extrapolation to sparsely supported treatment regions, where limited overlap can yield unstable estimates \citep{petersen2012diagnosing, kennedy2017non}, we compute MISE and ADRF error over the range between the 5th and 95th percentiles of the observed treatment distribution, following standard approaches to limited overlap \cite{crump2009dealing}. 

\subsection{Evaluation on Fully Synthetic Benchmark}
\label{Exp: fully synthetic}
%We first evaluate the proposed model on \textit{fully synthetic} datasets, where the complete data-generating %mechanism---including covariates, treatment assignment, and outcome generation---is known. 
%Because the ground-truth causal effects are available in this setting, both treatment effect estimation and %variable selection performance can be evaluated quantitatively.

%\textbf{Data Generation}
%For each sample, covariates $X \in \mathbb{R}^{100}$ are drawn from a standard multivariate Gaussian %distribution. 
%We designate $\{X_1,\ldots,X_5\}$ as confounders, $\{X_6,\ldots,X_{10}\}$ as outcome predictors, $\{X_{11},\ldots,X_{15}\}$ as instrumental variables, and $\{X_{16},\ldots,X_{100}\}$ as noise variables.
%We define two coefficient vectors, $c_t$ and $c_o$, for treatment and outcome. 
%For each covariate, the corresponding coefficient is sampled from $\mathcal{U}(0.5, 1.0)$ if the covariate %affects the treatment or outcome and is set to zero otherwise. 
%Thus, confounders have nonzero coefficients in both $c^{t}$ and $c^{o}$, outcome predictors only in $c^{o}$, %instruments only in $c^{t}$, and noise variables in neither. 

%The treatment is generated by min-max normalizing
%\begin{align*}
%a_{i}^* = \sum_{j=1}^{100} c_j^t X_{ij} + \epsilon_i^{A}, \quad \epsilon_i^{A} \sim \mathcal{N}(0, 0.1^2).
%\end{align*}
%as 
%\begin{align*}
%a_{i} = \frac{a_{i}^* - \min_{k}(a_{k}^*)}{\max_{k}(a_{k}^*) - \min_{k}(a_{k}^*)}.
%\end{align*}
%The potential outcome at treatment level $a$ is generated using a non-linear dose-response function:
%\begin{align*}
%Y_{i}(a) = \sum_{j=1}^{100} c_j^o X_{ij} + 4 \left( \frac{a^2}{0.8^2 + a^2} \right) + \epsilon_{i}^{Y},
%\end{align*}
%where $\epsilon_{i}^{Y} \sim \mathcal{N}(0, 0.1^2)$.
%Here, the first term represents the baseline outcomes from the outcome-relevant covariates, while the second term %specifies a nonlinear saturating Hill-type dose-response curve scaled by $4$.
%For evaluation, we generate datasets with $N=2000$ samples and repeat the experiment over $10$ runs.

\textbf{Data Generation.}
For each replicate, we generate \(N=2000\) observations with \(\mathbf{x}_i\sim\mathcal{N}(\mathbf{0},I_{100})\).
We assign \(X_{1:5}\) to the confounder role, \(X_{6:10}\) to the outcome-predictor role, \(X_{11:15}\) to the instrumental-variable role, and \(X_{16:100}\) to the irrelevant role. 

Let \(\mathbf{c}^{t},\mathbf{c}^{o}\in\mathbb{R}^{100}\) denote the treatment and outcome coefficient vector. 
Their active entries are drawn once from \(\mathcal{U}(0.5,1.0)\) and held fixed across all replicates; each replicate independently draws a new covariate matrix and noise terms.
Confounders are active in both vectors, outcome predictors only in \(\mathbf{c}^{o}\), instruments only in \(\mathbf{c}^{t}\), and irrelevant variables in neither.

We first generate the latent treatment score as
\[
a_i^{*}
=
(\mathbf{c}^{t})^{\top}\mathbf{x}_i+\epsilon_i^{A},
\qquad
\epsilon_i^{A}\sim\mathcal{N}(0,0.1^2),
\]
and rescaled it to \([0,1]\) within each replicate:
\[
a_i=\frac{a_i^{*}-\min_k a_k^{*}}{\max_k a_k^{*}-\min_k a_k^{*}}.
\]

The ground-truth conditional dose-response function is 
\[
\mu(a,\mathbf{x}_i)=(\mathbf{c}^{o})^{\top}\mathbf{x}_i+4(\frac{a^2}{0.8^2+a^2}),
\]
and the observed outcome is generated as
\[
y_i=\mu(a_i,\mathbf{x}_i)+\epsilon_i^{Y}, \quad \epsilon_i^{Y}\sim\mathcal{N}(0,0.1^2).
\]
The first term captures covariate-dependent baseline outcomes, while the second specifies a nonlinear saturating treatment-response component. 
%We report the mean and standard deviation over 10 independently generated replicates. 


\begin{table*}
\centering
\caption{\textbf{(RQ1)} 
Treatment effect estimation on the fully synthetic benchmark. 
Results are reported as mean \(\pm\) standard deviation over 10 independently generated replicates. Best results are shown in bold; lower is better.}%MISE represents the Mean Integrated Squared Error, measuring individual-level estimation accuracy, while $\epsilon_{\text{ADRF}}$ denotes the population-level Average Dose-Response Function error.}
\footnotesize
\label{Table1}
\begin{tabular}{lcccc}
\toprule
 & \multicolumn{2}{c}{In-sample} & \multicolumn{2}{c}{Out-of-sample} \\
\cmidrule(lr){2-3} \cmidrule(lr){4-5}
Model & MISE & $\epsilon_{\text{ADRF}}$ & MISE & $\epsilon_{\text{ADRF}}$ \\
\midrule
DRNet & 1.6278 $\pm$ 1.6570 & 0.4082 $\pm$ 0.2250 & 1.1689 $\pm$ 1.8786 & 0.3908 $\pm$ 0.2498 \\
SCIGAN & 0.4284 $\pm$ 0.0190 & 0.2968 $\pm$ 0.0123 & 0.4289 $\pm$ 0.0267 & 0.2971 $\pm$ 0.0116 \\
VCNet & 0.7366 $\pm$ 0.0471 & 0.4954 $\pm$ 0.0347 & 0.7400 $\pm$ 0.0382 & 0.4965 $\pm$ 0.0372 \\
IBEX & 0.8028 $\pm$ 0.0498 & 0.5562 $\pm$ 0.0391 & 0.8101 $\pm$ 0.0447 & 0.5545 $\pm$ 0.0392 \\
GIKS & 0.4857 $\pm$ 0.0894 & 0.3233 $\pm$ 0.0625 & 0.4946 $\pm$ 0.0886 & 0.3239 $\pm$ 0.0604 \\
DDMLCT & 1.6900 $\pm$ 0.0196 & 0.4917 $\pm$ 0.0287 & 1.6798 $\pm$ 0.0241 & 0.5099 $\pm$ 0.0260 \\
ACFR & 0.8937 $\pm$ 0.2888 & 0.5473 $\pm$ 0.0775 & 0.9027 $\pm$ 0.3153 & 0.5502 $\pm$ 0.0812 \\
\rowcolor{blue!10}
\textbf{CVSCEN} & \textbf{0.0831 $\pm$ 0.0148}& \textbf{0.0152$\pm$ 0.0107}& \textbf{0.0860$\pm$ 0.0142}& \textbf{0.0148$\pm$ 0.0109}\\
\bottomrule
\end{tabular}
\end{table*}

\begin{table}[t!]
\centering
\caption{\textbf{(RQ2)} Variable selection on the fully synthetic benchmark.  
We report role-specific FDR and TPR for confounders $c$, outcome predictors $p$, and their union $c \cup p$. 
Results are shown as mean $\pm$ standard deviations over 10 independently generated replicates. }
\label{Table2}
\resizebox{\columnwidth}{!}{
\begin{tabular}{lccc}
\toprule
Metric & $c$ & $p$ & $c \cup p$ \\
\midrule
FDR & 0.0542$\pm$ 0.1179& 0.0000$\pm$ 0.0000& 0.0231$\pm$ 0.0692
\\
TPR & 1.0000$\pm$ 0.0000& 0.9800$\pm$ 0.0600& 1.0000 $\pm$ 0.0000 \\
\bottomrule
\end{tabular}
}

\vspace{1ex}
\end{table}



\textbf{Results}
\cref{Table1} and \cref{Table2} show that CVSCEN achieves the strongest overall performance on the fully synthetic benchmark, combining accurate dose-response estimation with effective recovery of causally relevant variables. 
For treatment effect estimation, CVSCEN obtains the lowest errors among all evaluated methods. In the out-of-sample setting, it achieves an MISE of 0.0794 and an ADRF error of 0.0166, compared with an MISE of 0.4289 for the strongest baseline, SCIGAN.

%This outstanding performance is primarily driven by CVSCEN's remarkable variable selection capabilities. As shown in \cref{Table2}, CVSCEN successfully isolates all true confounders without missing a single one ($c$-TPR 1.0000) while maintaining a nearly zero False Discovery Rate ($c$-FDR 0.0167). Similarly, it recovers almost all outcome predictors ($p$-TPR 0.9800) with absolutely zero false discoveries ($p$-FDR 0.0000). Overall, for the union of causally relevant variables ($c \cup p$), it achieves a perfect True Positive Rate of 1.0000 and a perfect False Discovery Rate of 0.0000. 
For variable selection, CVSCEN recovers all confounders and nearly all outcome predictors, with a TPR of 1.0000 and 0.9800, respectively. 
For their union, it achieves a TPR of 1.0000 and a low FDR of 0.0231, indicating that it retains the variables relevant for adjustment or outcome modeling while largely excluding instrumental and irrelevant variables. These results are consistent with the intended benefit of explicit role-aware selection in reducing high-dimensional nuisance information.
%This near-flawless feature selection enables the model to completely filter out the overwhelming number of purely noisy variables and bias-amplifying instrumental variables present in the high-dimensional synthetic dataset, which explains why CVSCEN maintains such exceptional individual-level and population-level prediction accuracy compared to baseline methods that indiscriminately process all observed variables.

\subsection{Evaluation on Semi-Synthetic IHDP Benchmark}
\label{Exp: semi-synthetic}
We next evaluate CVSCEN on a semi-synthetic benchmark based on the Infant Health and Development Program (IHDP) dataset that combines real-world covariates with known treatment and outcome mechanisms. 

\textbf{Data Generation.}
We use the 25-dimensional covariates from the IHDP dataset and generate continuous treatments and outcomes following the nonlinear data-generating process used in VCNet \citep{nie2021vcnet}. 
The covariates are partitioned into confounders $c = \{1,2,3,5,6\}$, outcome predictors $p =\{4, 7, \ldots, 15\}$, and instrumental variables $w  =\{16,\ldots,25\}$.

The continuous treatment $a$ is generated using a sigmoid transformation $a = \sigma(2 l_{a})$,
%\begin{equation}
    %a = \frac{1}{1 + \exp(-2 \cdot \ell_a)},
%\end{equation}
where $\ell_a$ is defined as a nonlinear function of confounders and instrumental variables:
\begin{multline}
\ell_a = \frac{x_1}{1 + x_2} + \frac{\max(x_3, x_5, x_6)}{0.2 + \min(x_3, x_5, x_6)} \\
+ \tanh(5 \cdot (\bar{X}_w - \mu_w)) - 2.0 + \eta.
\end{multline}
Here, $\bar{X}_{w}$ denotes the mean of the instrumental variables, $\mu_{w}$ is its population mean, and $\eta \sim \mathcal{N}(0, 0.5^2)$ represents the treatment assignment noise.


The outcome $Y$ is generated as:
\begin{multline}
Y = \frac{\sin(3\pi a)}{1.2 - a} \Bigg( 1.5 \cdot \tanh(5 \cdot (\bar{X}_p - \mu_p)) \\
+ \frac{0.5 \cdot \exp(0.2(x_1 - x_6))}{0.1 + \min(x_2, x_3, x_5)} \Bigg) + \epsilon
\end{multline}
where $\bar{X}_{p}$ is the mean of the outcome predictors, $\mu_{p}$ is their population mean, and $\epsilon \sim \mathcal{N}(0, 0.5^2)$. 
This benchmark combines nonlinear treatment assignment, nonlinear outcome generation, instrumental variables, and naturally correlated real-world covariates.
%making it suitable for evaluating both estimation accuracy and variable selection. 

\begin{table*}
\centering
\caption{\textbf{(RQ1)} Treatment effect estimation on the semi-synthetic IHDP benchmark.}
\footnotesize
\label{Table3}
\begin{tabular}{lcccc}
\toprule
 & \multicolumn{2}{c}{In-sample} & \multicolumn{2}{c}{Out-of-sample} \\
\cmidrule(lr){2-3} \cmidrule(lr){4-5}
Model & MISE & $\epsilon_{\text{ADRF}}$ & MISE & $\epsilon_{\text{ADRF}}$ \\
\midrule

DRNet & 1.1971 $\pm$ 0.0604 & 0.5770 $\pm$ 0.0995 & 1.2481 $\pm$ 0.0745 & 0.6047 $\pm$ 0.1005 \\
SCIGAN & 2.4696 $\pm$ 0.1497 & 1.6439 $\pm$ 0.1000 & 2.4955 $\pm$ 0.1855 & 1.6855 $\pm$ 0.0796 \\
VCNet & 0.8188 $\pm$ 0.1509 & 0.2021 $\pm$ 0.0575 & 0.8455 $\pm$ 0.1737 & 0.2270 $\pm$ 0.0609 \\
IBEX & 0.5507 $\pm$ 0.0618& 0.1628 $\pm$ 0.0582& 0.6182 $\pm$ 0.1176 & 0.1929 $\pm$ 0.0620 \\
GIKS & 0.7755 $\pm$ 0.1213 & 0.2514 $\pm$ 0.1035 & 0.7358 $\pm$ 0.1177 & 0.3008 $\pm$ 0.1076 \\
DDMLCT & 4.4405 $\pm$ 3.5042 & 4.1057 $\pm$ 3.6938 & 4.4710 $\pm$ 3.4871 & 4.1133 $\pm$ 3.6971 \\
ACFR & 1.2622 $\pm$ 0.1216 & 0.4857 $\pm$ 0.1176 & 1.2723 $\pm$ 0.1900 & 0.4998 $\pm$ 0.1609 \\
\rowcolor{blue!10}
\textbf{CVSCEN} & \textbf{0.5181 $\pm$ 0.0444}& \textbf{0.1332 $\pm$ 0.0281}& \textbf{0.5316 $\pm$ 0.0774}& \textbf{0.1463 $\pm$ 0.0405}\\
\bottomrule
\end{tabular}
\end{table*}

\begin{table}[h]
\centering
\caption{\textbf{(RQ2)} Variable selection on the semi-synthetic IHDP benchmark. We report FDR and TPR for confounders $c$, outcome predictors $p$, and their union $c \cup p$. }
\label{Table4}
\resizebox{\columnwidth}{!}{
\begin{tabular}{lccc}
\toprule
Metric & $c$ & $p$ & $c \cup p$ \\
\midrule
FDR & 0.0833 $\pm$ 0.0833& 0.0000 $\pm$ 0.0000& 0.0000 $\pm$ 0.0000
\\
TPR & 1.0000 $\pm$ 0.0000 & 0.5100 $\pm$ 0.0300& 0.7067 $\pm$ 0.0442\\
\bottomrule
\end{tabular}
}
\end{table}

\textbf{Results.} 
\cref{Table3} and \cref{Table4} show that CVSCEN achieves the best overall performance on the semi-synthetic IHDP benchmark while providing explicit covariate-level selection. 
%CVSCEN achieves the best performance across all estimation metrics.
%surpassing all evaluated methods.
It obtains the lowest out-of-sample MISE of 0.5316 and ADRF error of 0.1463, as well as the lowest in-sample MISE of 0.5181 and ADRF error of 0.1332.
These results indicate that CVSCEN maintains strong estimation accuracy under realistic covariate correlations. 
%CVSCEN not only remains highly competitive but actually 
%outperforms state-of-the-art methods like IBEX under realistic covariate correlations, all while providing explicit and interpretable variable selection.

For variable selection, CVSCEN retains a substantial subset of the causally relevant variables while excluding nuisance variables. 
For $c \cup p$, it achieves an FDR of 0.0000 and a TPR of 0.7067, indicating that none of the selected variables are nuisance variables and that the model retains many covariates relevant for adjustment or outcome modeling. 
Role-specific results further show that CVSCEN recovers all confounders in this benchmark, with a confounder TPR of 1.0000. 
For outcome predictors, it achieves an FDR of 0.0000 and a lower TPR of 0.5100, indicating more conservative selection of prognostic variables. 
%This conservative selection may reflect the strong correlations among the IHDP covariates: when several outcome predictors contain redundant prognostic information, the outcome model can achieve accurate prediction using only a subset of them. 
%Accordingly, the lower outcome-predictor TLR does not necessarily imply a proportional loss in treatment effect estimatoin accuracy. 
This may reflect redundancy among the correlated IHDP covariates, as the outcome model can obtain sufficient prognostic information from only a subset of highly correlated predictors. 
Accordingly, incomplete recovery of all outcome predictors need not translate into a proportional loss in treatment effect estimation accuracy.
%exhibit strong multicollinearity. While the independence regularization successfully prevents false discoveries, the model can sufficiently predict the outcome using only a subset of these highly correlated outcome predictors, leading it to conservatively select fewer variables.

%Compared with baselines that use all observed covariates as input, CVSCEN can reduce the influence of instrumental and irrelevant variables through explicit selection. 
%IBEX also performs strongly, likely owing to its representation bottleneck, but it does not provide covariate-level role assignments. 
%DDMLCT exhibits substantially larger errors and higher variability, suggesting sensitivity to nuisance estimation under a limited sample size of this benchmark. 
%Overall, an IHDP results show that CVSCEN can combine accurate continuous treatment effect estimation with sparse and interpretable variable selection under realistic covariate correlations.  
%in semi-synthetic settings with realistic covariates.


\subsection{Ablation Study}
\label{sec:ablation}
We extend our ablation study to validate not only the internal mechanisms of CVSCEN but also the necessity of its core architectural designs: explicit variable selection and the specialized outcome model. We evaluate five variants of CVSCEN on the semi-synthetic IHDP benchmark: (1) \textbf{w/o Variable Selection}, which feeds all covariates into the outcome model without selection; (2) \textbf{w/ Simple MLP}, which retains variable selection but replaces the dual-stream outcome model with a standard MLP; (3) \textbf{Oracle}, which provides the ground-truth confounders and outcome predictors to the outcome model; (4) \textbf{w/o HSIC Loss}, which removes the treatment-independence regularizer; and (5) \textbf{w/o Correlation Prior}, which omits the dependence-based guidance. The latter two isolate the contribution of individual selection mechanisms. All other hyperparameters and training procedures are held fixed.

\begin{figure}[t]
    \centering
    \includegraphics[width=0.48\textwidth]{figures/ablation_ihdp.png}
    \caption{\textbf{(RQ3)} Treatment effect estimation error and variable selection performance of the models in the ablation study. (a) MISE, (b) ADRF error, (c) FDR, (d) TPR.}
    \label{fig:ablation_ihdp}
\end{figure}

\textbf{Necessity and Efficacy of Variable Selection.} 
Comparing CVSCEN with the \textit{w/o Variable Selection} variant highlights the critical importance of filtering covariates. Indiscriminately incorporating all covariates yields the highest estimation error (MISE 1.3934), as the model is forced to condition on non-prognostic instrumental and irrelevant variables, severely destabilizing the continuous dose-response estimation. Conversely, the full CVSCEN model achieves an MISE of 0.5482, closely matching the \textit{Oracle} performance (MISE 0.5419). This indicates that CVSCEN successfully filters out nuisance variables and identifies the necessary causal roles almost as effectively as if the ground-truth roles were known a priori.

\textbf{Importance of the Specialized Outcome Model.} 
Replacing our dual-stream architecture with a standard prediction network (\textit{w/ Simple MLP}) substantially degrades performance (MISE 1.2467). Crucially, this simplified outcome model not only increases prediction error but also devastates the variable selection itself: the True Positive Rate for outcome predictors (TPR for $p$) plummets to near zero. Because a simple MLP struggles to separate the treatment signal from high-dimensional prognostic variations, it cannot properly utilize outcome predictors. Consequently, the end-to-end variable selector receives no useful feedback to identify them. This highlights a key synergy in CVSCEN: a robust, specialized outcome model is required to provide the stable gradient signals necessary for learning accurate role assignments.

\textbf{Role of Individual Selection Mechanisms.} 
Finally, our internal ablations confirm the complementary roles of the selection regularizers. Removing the dependence-based guidance (\textit{w/o Correlation Prior}) degrades role separation, causing confounder FDR to surge as outcome predictors are misclassified. Meanwhile, removing the treatment-independence regularizer (\textit{w/o HSIC Loss}) increases both estimation error (MISE 0.6573) and variance. Without HSIC, treatment-associated covariates can leak into the outcome-predictor branch, entangling the prognostic representations with the treatment assignment. Together, these mechanisms ensure both correct variable assignment and stable dose-response estimation.

\section{Conclusion}
We introduced CVSCEN, an end-to-end framework for continuous treatment effect estimation with explicit causal role-aware variable selection. 
CVSCEN jointly learns confounder, outcome-predictor, and excluded roles with a dual-stream outcome model, combining one-sided dependence guidance, HSIC-based regularization, and role-specific sparsity. 
Experiments on synthetic and semi-synthetic benchmarks show improved dose-response estimation while recovering covariates relevant to adjustment and outcome modeling.

CVSCEN assumes an observed sufficient adjustment set, does not address unobserved confounding or severe overlap violations, and is evaluated only where causal roles are known. Future work will study real-world applications, more complex treatment settings, and theoretical guarantees for role recovery and dose-response estimation.
}

\newpage
\bibliography{cvscen.bib}
\bibliographystyle{icml2026}

\newpage
\onecolumn
\definecolor{darkblue}{rgb}{0.0, 0.1, 0.5}
\newpage
\appendix

\begin{center}
  {\fontsize{16pt}{18pt}\selectfont\bfseries\color{darkblue} Appendix for the Paper} \\[0.6em]
  {\fontsize{14pt}{16pt}\selectfont \color{darkblue} ``Adaptive Variable Selection for Continuous Treatment Effect Estimation'} \\[1.5em]
\end{center}
{\color{darkblue}
\textbf{A \quad Notations}  
\vspace{-0.5em}  

\textbf{B \quad Algorithms}  
\vspace{0.5em}  

\textbf{C \quad Additional Explanations about the Experiments}  
\vspace{-0.5em}  
\begin{itemize}[leftmargin=1.5em, itemsep=0em]
    \item C.1 \quad Datasets
    \item C.2 \quad Baselines
    \item C.3 \quad Implementationl Details
\end{itemize}

\textbf{D \quad Additional Experimental Results} 
\vspace{0.5em}  

\textbf{E \quad Proofs of the Theorems}  
\vspace{-0.5em}  
\begin{itemize}[leftmargin=1.5em, itemsep=0em]
    \item E.1 \quad Proof of Thm 1
    \item E.2 \quad Proof of Thm 2
\end{itemize}}
\newpage

\section{Notations}
\begin{table}[h!]
\centering
\begin{tabular}{|c|l|}
\hline
\textbf{Notation} & \textbf{Description} \\ \hline
\( A \) & Continuous treatment variable. \\ \hline
\( Y \) & Outcome variable. \\ \hline
\( X = (X_1, X_2, \dots, X_d) \) & Covariate vector (with \( d \) dimensions). \\ \hline
\( Y_i(a) \) & Potential outcome for unit \( i \) under treatment level \( a \). \\ \hline
\( \text{DRF}_i(a) \) & Individual-level dose-response function. \\ \hline
\( \text{CADRF}(a, x) = \mathbb{E}[Y(a) | X = x] \) & Conditional average dose-response function. \\ \hline
\( \text{ADRF}(a) \) & Average dose-response function at the population level. \\ \hline
\( \ell_{j,r} \) & Logit (assignment score) for covariate \( X_j \) being assigned to role \( r \in \{0, C, P\} \). \\ \hline
\( z_{j,r}^{(e)} \) & Relaxed assignment probability for covariate \( X_j \) to role \( r \) at epoch \( e \). \\ \hline
\( \rho_j \) & HSIC between \( X_j \) and treatment \( A \),\\ \hline
%used to guide the selection of confounders. \\ \hline
\end{tabular}
\caption{Notations and Key Terms}
\end{table}

\section{Algorithms}

\section{Additional Explanations about the Experiments}
\subsection{IHDP Datasets.}
The Infant Health and Development Program (IHDP) dataset is widely utilized as a standard benchmark in causal inference to evaluate treatment effect estimation algorithms. Originally collected from a randomized study targeting low-birth-weight and premature infants, the dataset provides a rich set of real-world covariates. 

Specifically, we use 25 covariates for each subject, encompassing both child characteristics (e.g., birth weight, head circumference) and maternal attributes (e.g., maternal age, education level, and behaviors during pregnancy). Among the 25 covariates, 6 are continuous, and 19 are binary. To ensure stable model training, the continuous covariates are standardized to have zero mean and unit variance, while the binary covariates are retained in their original $\{0, 1\}$ format.

While the IHDP dataset is traditionally used for binary treatment settings, we adapt it for our continuous treatment semi-synthetic experiment. We detach the original treatment and outcome variables, retaining only the 25-dimensional real-world covariate matrix $X \in \mathbb{R}^{25}$. These covariates exhibit inherent, natural correlation structures and complex multicollinearity. By systematically partitioning these real-world covariates into distinct causal roles—confounders, outcome predictors, and instrumental variables—we construct a highly challenging and realistic environment to validate the model's capacity to estimate simultaneous continuous treatment effects and select causal variables.


\subsection{Baselines.}
We compare the proposed method with representative baselines.
\begin{itemize}
    \item Base (Mean Baseline): A naive baseline that estimates the average outcome of the population.
    \item VCNet (Variable Coefficient Network) \citep{nie2021vcnet}: A model for continuous treatments that uses varying coefficient layers to approximate the dose-response function. 
    \item DRNet (Dose Response Network) \citep{schwab2020learning}: A deep learning architecture designed for continuous treatments that partitions the dosage range into multiple strata and learns separate head networks to estimate the individual dose-response function.
    \item SCIGAN \citep{bica2020estimating}: A generative framework that utilizes Generative Adversarial Networks (GANs) to predict counterfactual outcomes for continuous-valued interventions. 
    \item GIKS (Gradient Interpolation and Kernel Smoothing) \citep{nagalapatti2024continuous}: A recent method that leverages gradient interpolation and kernel smoothing to robustly estimate continuous treatment effects.
    \item IBEX (Cauchy-Schwarz Bottleneck) \citep{van2026continuous}: A representation learning approach that uses an Information Bottleneck based on Cauchy-Schwarz divergence to compress non-causal variables.
    \item DDMLCT (Double Machine Learning for Continuous Treatments) \citep{colangelo2025double}: A double machine learning framework designed for continuous treatments with high-dimensional nuisance parameters.
    \item ACFR (Adversarial Counterfactual Regression) \citep{kazemi2024adversarially}: A deep learning method that utilizes adversarially balanced representations to mitigate confounding bias in continuous treatment effect estimation.
    \item CVSCEN (Ours): Our proposed end-to-end framework that simultaneously performs causal role-aware variable selection and continuous treatment effect estimation. 
\end{itemize}

\subsection{Implementation Details}

\textbf{Network Architecture.} The CVSCEN framework is implemented as an end-to-end neural network consisting of two primary functional modules.

\begin{itemize}
    \item \textbf{Dynamic Variable Selection Layer}  This layer maintains a learnable logit matrix of size $d \times 3$, where $d$ is the number of input covariates. During training, it produces continuous masks $M \in [0, 1]^d$ using the Gumbel-Softmax trick to ensure stable gradient flow, while switching to discrete one-hot selections $M \in \{0, 1\}^d$ during test via argmax. The Gumbel-Softmax temperature $\tau$ anneals from 10.0 to 0.1 over the training epochs. HSIC regularization is further applied to penalize dependencies between selected variables and the treatment, where applicable.
    
    \item \textbf{Outcome Prediction Layer} receives the covariates masked as confounders ($x_c$) and outcome predictors ($x_p$), alongside the 1-dimensional observed continuous treatment. Its internal architecture utilizes a dual-stream feature extraction strategy: both $x_c$ and $x_p$ are independently processed by deep feature extractors containing Residual Blocks and LayerNorm, utilizing SiLU activations. The continuous treatment is expanded into a higher-dimensional space using NeRF-style positional encoding (e.g., $a, \sin(2^k \pi a), \cos(2^k \pi a)$). The extracted confounder representations are then dynamically modulated via a treatment-conditioned affine transformation (generating scale $\gamma$ and shift $\beta$). The prognostic representations are kept strictly independent of the treatment. Finally, the modulated representations are aggregated and passed through a prediction MLP to output the scalar response.
\end{itemize}

\textbf{Training and Hyperparameter Tuning Details.} 
All experiments were implemented using the PyTorch framework. To ensure a fair comparison and optimal performance, we systematically tuned the hyperparameters---including the learning rate and weight decay---for both our proposed CVSCEN model and all baseline methods. For CVSCEN, the network parameters were optimized using the Adam optimizer. To ensure stable convergence throughout training, we coupled Adam with a StepLR scheduler that decays the learning rate by a factor of $\gamma=0.97$ every 100 epochs, and applied $L_2$ regularization (weight decay) to prevent overfitting and control the complexity of the deep neural network.


For robust evaluation, the datasets were randomly partitioned into training, validation, and test sets using a $63{:}27{:}10$ split ratio. The training set was used to update the model weights, while the validation set was utilized exclusively for hyperparameter optimization via Optuna-based Bayesian optimization using the Tree-structured Parzen Estimator (TPE) algorithm to efficiently explore the search space. Because our benchmarks provide ground-truth dose-response functions, we directly minimized the out-of-sample Mean Integrated Squared Error (MISE) on the validation set during tuning. 

The search spaces were carefully selected to cover standard configurations and authors' recommendations for continuous treatment effect estimation. The comprehensive search ranges for CVSCEN and the baseline models are summarized in Table \ref{tab:hyper_appendix}. The final performance of the causal effect estimation was evaluated exclusively on the held-out test set.


\begin{table}[htbp]
\centering
\caption{Hyperparameter Grid Search Spaces}
\label{tab:hyper_appendix}
\resizebox{\columnwidth}{!}{
\begin{tabular}{ll}
\toprule
Model & Search Range \\
\midrule
\multirow{5}{*}{\textbf{Common}} & Learning Rate: $\{10^{-4}, 5\times 10^{-4}, 10^{-3}, 5\times 10^{-3}, 10^{-2}\}$ \\
& Hidden Dimension: $\{32, 64, 128, 256\}$ \\
& Epochs: $\{300, 400, 500\}$ \\
& Batch Size: $\{32, 64, 128\}$ \\
& Weight Decay: $\{10^{-4}, 10^{-3}, 5\times 10^{-3}, 10^{-2}\}$ \\
\midrule
\multirow{4}{*}{\textbf{CVSCEN (Ours)}} & HSIC Penalty ($\lambda_{\text{HSIC}}$): $\{0.5, 1.0, 2.5, 5.0, 10.0\}$ \\
& Dependence Guidance Weight ($\beta$): $\{1.0, 10.0, 50.0, 100.0, 500.0\}$ \\
& Sparsity Penalty ($\lambda_{\mathrm{reg},C}, \lambda_{\mathrm{reg},P}$): $\{0.01, 0.05, 0.1, 0.25, 0.5\}$ \\
& Initial Confounder Logit ($\ell_{C, init}$): $\{0.1, 1.0, 10.0\}$ \\
\midrule
\textbf{VCNet, DRNet} & $\alpha$: $\{0.1, 0.5, 1.0, 2.0\}$, $n_{\text{grid}}$: $\{5, 10, 20\}$, degree: $\{2, 3\}$, Epochs: $\{300, 500, 700, 800\}$, TR $lr$: $\{10^{-4}, 10^{-3}, 10^{-2}\}$ \\
\midrule
\textbf{ACFR} & $lr_s$: $\{10^{-5}, 10^{-4}, 10^{-3}\}$, $\gamma_1$: $\{0.1, 1.0, 5.0, 10.0\}$, $\gamma_2$: $\{0.1, 0.2, 0.5\}$, $std$: $\{0.1, 0.2, 0.5\}$ \\
\midrule
\textbf{SCIGAN} & Batch Size: $\{16, 32, 64, 128, 256\}$, $\alpha$: $\{0.1, 1.0, 10.0\}$, $n_{\text{samples}}$: $\{3, 5, 10\}$, $lr_g, lr_d$: $\{10^{-4}, 5\times 10^{-4}, 10^{-3}\}$, $d_{\text{steps}}$: $\{1, 3, 5\}$ \\
\midrule
\textbf{GIKS} & Epochs: $\{200, 300, 400, 500\}$, $\lambda_{\text{GIKS}}$: $\{10^{-5}, 10^{-4}, 10^{-3}, 10^{-2}\}$, $\delta$: $\{0.01, 0.05, 0.1\}$, $n_{\text{grid}}$: $\{5, 10, 20\}$, degree: $\{2, 3\}$ \\
\midrule
\multirow{2}{*}{\textbf{IBEX}} & $\beta$: $\{0.001, 0.01, 0.1\}$, $\gamma$: $\{0.001, 0.1, 1.0, 10.0\}$, $z_{\text{dim}}$: $\{16, 32, 64\}$, $t_{\text{dim}}$: $\{4, 8, 16\}$ \\
& Epochs: $\{1000, 2000, 3000\}$, Hidden Dimension: $\{128, 256, 512\}$ \\
\midrule
\multirow{2}{*}{\textbf{DDMLCT}} & $lr_s$: $\{0.001, 0.01, 0.05, 0.15, 0.4\}$, $L$: $\{2, 5, 10\}$, $lr$: $\{10^{-4}, 5\times 10^{-4}, 10^{-3}, 5\times 10^{-3}, 10^{-2}, 0.15, 0.4\}$ \\
& Hidden Dimension: $\{10, 25, 64, 100, 128\}$, Weight Decay: $\{10^{-4}, 10^{-3}, 5\times 10^{-3}, 10^{-2}, 0.1, 0.2, 0.3\}$ \\
\bottomrule
\end{tabular}
}
\end{table}

\section{Additional Experimental Results}
\begin{figure*}[t]
    \centering

    \begin{minipage}[b]{0.24\textwidth}
        \centering
        \includegraphics[width=\textwidth]{figures/drnet_ihdp.png}
    \end{minipage}
    \hfill
    \begin{minipage}[b]{0.24\textwidth}
        \centering
        \includegraphics[width=\textwidth]{figures/scigan_ihdp.png}
    \end{minipage}
    \hfill
    \begin{minipage}[b]{0.24\textwidth}
        \centering
        \includegraphics[width=\textwidth]{figures/vcnet_ihdp.png}
    \end{minipage}
    \hfill
    \begin{minipage}[b]{0.24\textwidth}
        \centering
        \includegraphics[width=\textwidth]{figures/ibex_ihdp.png}
    \end{minipage}

    \vspace{0.5em}

    % Row 2
    \begin{minipage}[b]{0.24\textwidth}
        \centering
        \includegraphics[width=\textwidth]{figures/giks_ihdp.png}
    \end{minipage}
    \hfill
    \begin{minipage}[b]{0.24\textwidth}
        \centering
        \includegraphics[width=\textwidth]{figures/ddmlct_ihdp.png}
    \end{minipage}
    \hfill
    \begin{minipage}[b]{0.24\textwidth}
        \centering
        \includegraphics[width=\textwidth]{figures/acfr_ihdp.png}
    \end{minipage}
    \hfill
    \begin{minipage}[b]{0.24\textwidth}
        \centering
        \includegraphics[width=\textwidth]{figures/cvscen_ihdp.png}
    \end{minipage}

    \caption{Average Dose-Response Function (ADRF) estimated by each model on the IHDP dataset. The respective model name is indicated in the title of each subplot.}
    \label{fig:adrf_all}
\end{figure*}

\begin{figure*}[t]
    \centering

    % Row 1
    \begin{minipage}[b]{0.24\textwidth}
        \centering
        \includegraphics[width=\textwidth]{figures/drnet_synt.png}
    \end{minipage}
    \hfill
    \begin{minipage}[b]{0.24\textwidth}
        \centering
        \includegraphics[width=\textwidth]{figures/scigan_synt.png}
    \end{minipage}
    \hfill
    \begin{minipage}[b]{0.24\textwidth}
        \centering
        \includegraphics[width=\textwidth]{figures/vcnet_synt.png}
    \end{minipage}
    \hfill
    \begin{minipage}[b]{0.24\textwidth}
        \centering
        \includegraphics[width=\textwidth]{figures/ibex_synt.png}
    \end{minipage}

    \vspace{0.5em}

    \begin{minipage}[b]{0.24\textwidth}
        \centering
        \includegraphics[width=\textwidth]{figures/giks_synt.png}
    \end{minipage}
    \hfill
    \begin{minipage}[b]{0.24\textwidth}
        \centering
        \includegraphics[width=\textwidth]{figures/ddmlct_synt.png}
    \end{minipage}
    \hfill
    \begin{minipage}[b]{0.24\textwidth}
        \centering
        \includegraphics[width=\textwidth]{figures/acfr_synt.png}
    \end{minipage}
    \hfill
    \begin{minipage}[b]{0.24\textwidth}
        \centering
        \includegraphics[width=\textwidth]{figures/cvscen_synt.png}
    \end{minipage}

    \caption{Average Dose-Response Function (ADRF) estimated by each model on the synthetic dataset.
    The dashed black line denotes the true ADRF, and the solid red line denotes the predicted ADRF.
    The shaded region indicates the gap between the two. The respective model name is indicated in the title of each subplot.}
    \label{fig:adrf_all_synt}
\end{figure*}
\section{Proofs for the Theorems}
\label{Proofs for the Theorems}
\end{document}
